\section{The chain argument without short-cycle deletions}
\label{sec:random-no-deletion-chain}

At girth proportional to the logarithm of the rank, the graph itself supplies the girth needed by the parallel-link argument.
No packing estimate depending on the size of the forest excess is needed.
All logarithms have base two.

\begin{lemma}[Normalizing a cut between protected cyclic cores]
\label{lem:random-protected-core-cut}
Let $G$ be connected and cubic.
Let $A,B$ be disjoint connected subgraphs of minimum degree at least two.
Suppose a vertex partition separating $A$ from $B$ has at most $q$ crossing edges.
Then there is such a partition with at most $q$ crossing edges whose two shores are connected normalized poles.
Their ranks are at least $\beta(A)$ and $\beta(B)$.
\end{lemma}

\begin{proof}
Choose a minimum edge cut among partitions with all of $A$ on one side and all of $B$ on the other.
Each shore is connected.
Indeed, a component not containing its protected connected subgraph could be moved to the other shore, strictly reducing the cut.

A protected vertex has at least two incidences within its protected subgraph, so at most one cut incidence.
If an unprotected vertex had at least two cut incidences, moving that vertex to the opposite shore would reduce the cut.
Thus every vertex has at most one cut incidence.
The shores have degree two at their cut endpoints and degree three elsewhere.
Their ranks are at least the ranks of their respective protected subgraphs.
\end{proof}

\begin{lemma}[A positive forest excess gives two large shores]
\label{lem:random-no-deletion-chain}
Fix $c>0$.
For all sufficiently large $r$, the following holds.
Let $G$ be connected and cubic of rank $r$ and girth at least $c\log r$.
Let $P$ be a connected marked set with $|P|\le3\log r$.
If
\[
 \Pr(\xi_G-P\text{ is a forest})>\frac12,
\]
then $G$ has a cut of order at most $2A\log r$, where $A=6e^{168}$, with connected normalized shores of ranks at least
$r^{\min\{c,1\}/10}$.
The threshold depends on $c$, but not on the positive amount of forest excess.
\end{lemma}

\begin{proof}
Replace $c$ by $\min\{c,1\}$ if necessary.
Put $x=\log r$, $z=\log x$, and use
\[
 L=\lfloor Ax\rfloor,\qquad
 D=\lfloor cx/2\rfloor,\qquad
 t=\lfloor z-2\log z\rfloor,\qquad
 s=\lfloor(D-1)/2\rfloor.
\]
The graph $J=G-P$ has girth greater than $D$.
Thus the empty family is a maximal packing of cycles of length at most $D$ outside $P$.
In the notation of the endpoint proof, its deletion set is $S=\varnothing$.

The graph $J$ has $v=2r+O(x)$ vertices and degree deficit at most $3|P|=O(x)$.
Hence $Ld_J/(2v)=o(1)$.
The even spectral trace sum over $6x\le\ell\le L$, together with
$L^2 2^{-s}=o(1)$, gives
\[
 \mu=\sum_{\substack{C\subseteq J\\|C|\le L}}2^{-|C|}
 \ge\left(\frac14-o(1)\right)\ln(A/6)>20.
\]

The parallel-link estimate applies to this empty packing.
Every pair of disjoint cycles of length at most $L$ has at most
\[
 1+\frac{2L}{t+1}+2^{t+1}+578(L/D)^2
 \le\frac{3Ax}{z}
\]
links for sufficiently large $r$.
The exceptional-pair error in the established refined second-moment estimate of
\texttt{polynomial-cycle-error.pdf} is therefore
\[
 O\!\left(2^{3Ax/z}L^3 2^{-s}\right)
 =r^{-c/4+o(1)}=o(1).
\]
Only the first conclusion of that lemma is used.
It requires forest probability greater than one half and gives
$W_2>\mu^2-(1+o(1))\mu$.

Apply the anchor-component, far-pair, and weighted maximal-disjointness argument from the established endpoint forest proof in
\texttt{polynomial-cycle-error.pdf}.
For clarity, its output is a cycle $C\subseteq J$ of length at most $L$, a neighbor $u$ of $C$, and vertex-disjoint cycles
$C_1,\ldots,C_T$ in the component $H_0$ of $G-V(C)$ containing $P$.
These cycles separate $u$ from $P$ in $H_0$ and are nested in that order.
The weighted partner mass is at least $8/L$, so
\[
 T\ge\frac{8\,2^s}{L^3}
 \ge\frac{2}{A^3}\frac{r^{c/4}}{x^3}.
\]
The proof of this output uses no inverse bound on $\mu$ and no positive lower bound on the forest excess.

Choose a middle index $i$.
Let $U_i$ be the $u$-component of $H_0-V(C_i)$.
It contains $C_1,\ldots,C_{i-1}$.
The complementary induced subgraph $H_0-U_i$ is connected: it contains the connected cycle $C_i$, and every other component of $H_0-V(C_i)$ has an edge to that cycle.
It contains $C_i,\ldots,C_T$.
Thus the ranks of these two connected subgraphs are at least $i-1$ and $T-i+1$, respectively.

The boundary of $U_i$ in $G$ has at most $2L$ edges.
At most $L$ end on $C_i$, since a vertex of that cycle has at most one incidence outside it.
The remaining boundary edges end on the removed anchor cycle $C$, and there are at most $L$ of those.

In each of the two connected induced subgraphs above, repeatedly delete degree-one vertices.
The remaining two-cores are connected, disjoint, and have minimum degree at least two.
Deleting leaves preserves every cycle and the rank.
For large $r$, both cores therefore have rank at least $T/3$.
They are separated by the cut around $U_i$.
Apply Lemma~\ref{lem:random-protected-core-cut}.
Its normalized shores have ranks at least $T/3\ge r^{c/10}$ for sufficiently large $r$.
\end{proof}

\begin{lemma}[A quarter bound for any small cut with two large shores]
\label{lem:random-general-large-shore-quarter}
Fix positive constants $c,K,\delta$.
Suppose connected cubic graphs have rank $r\to\infty$, girth at least $c\log r$, and a cut of order $k\le K\log r$ whose shores are connected normalized poles of ranks at least $r^\delta$.
Then
\[
 c(G)\le\left(\frac14+o_{c,K,\delta}(1)\right)2^r.
\]
The cut is not required to have minimum cyclic order.
\end{lemma}

\begin{proof}
Write the shore ranks as $h,j$, so $r=h+j+k-1$.
If $k=1$, no cycle crosses the cut and $c(G)\le2^h+2^j=o(2^r)$.
Assume henceforth that $k\ge2$.
Theorem~\ref{thm:random-all-terminal-sectors-girth} applies to the first shore:
$k\le(\log h)^2$ for large $r$, and its required girth is $O_K(\log\log r)$.
Each even terminal sector therefore has at most $(1/2+o(1))2^h$ forests.
Lemma~\ref{lem:random-growing-port-aggregate} gives total forest mass at most
$(1/2+o(1))2^{j+k-1}$ in the other shore.
Every crossing cycle has a unique pair of forest restrictions, so its count is bounded by the product of these two quantities.
The internal cycle counts are at most $2^h+2^j=o(2^r)$.
All errors are uniform under the fixed constants in the statement.
\end{proof}

\begin{theorem}[A high-girth edge-count dichotomy]
\label{thm:random-high-girth-edge-dichotomy}
Fix $c>0$.
For all sufficiently large ranks $r$, let $G$ be a connected loopless cubic graph of girth at least $c\log r$.
Let $W$ be a shortest cycle and let $e=uv$ be any edge of $W$.
If
\[
 P_{uv}(G-e)>2^{r-2},
\]
then
\[
 c(G)<2^{r-1}.
\]
Equivalently, if $c(G)\ge2^{r-1}$, every edge of a shortest cycle belongs to at most $2^{r-2}$ cycles.
\end{theorem}

\begin{proof}
The cubic Moore bound gives $|V(W)|\le3\log r$ for large $r$.
The usual two-state injection, now using the existing edge $e$, gives
\[
 \frac{P_{uv}(G-e)}{2^{r-1}}
 \le\Pr(\xi_G-V(W)\text{ is a forest}).
\]
Explicitly, if $A=W-e$, a terminal path $Q$ gives the two distinct even states
$Q\cup\{e\}$ and $Q\mathbin\triangle A$.
Their restrictions outside $V(W)$ are forests.
The hypothesis therefore gives a marked forest probability greater than one half.

Lemma~\ref{lem:random-no-deletion-chain} supplies a cut of order at most $2A\log r$ with two polynomial-rank normalized shores.
Lemma~\ref{lem:random-general-large-shore-quarter} bounds the cycle count by
$(1/4+o(1))2^r$, which is strictly below $2^{r-1}$.
\end{proof}

\begin{corollary}[No logarithmic-girth minimal counterexample to the sharp order]
\label{cor:random-high-girth-minimal-counterexample}
Fix $c,K>0$.
Suppose the bound
\[
 c(H)\le2^{s-1}+K s2^{s/2}
\]
holds for every connected graph of rank $s=r-1$.
Then, for all sufficiently large $r$, every connected loopless cubic graph $G$ of rank $r$ and girth at least $c\log r$ satisfies
\[
 c(G)\le2^{r-1}+K r2^{r/2}.
\]
\end{corollary}

\begin{proof}
If $c(G)<2^{r-1}$, there is nothing to prove.
Otherwise choose an edge $e$ of a shortest cycle.
The theorem gives $P_e(G-e)\le2^{r-2}$.
The graph $G-e$ is connected and has rank $r-1$, so the assumed predecessor bound yields
\[
 c(G)=c(G-e)+P_e(G-e)
 \le2^{r-1}+K(r-1)2^{(r-1)/2}
 <2^{r-1}+Kr2^{r/2}.
\]
This argument requires the induction hypothesis for general connected graphs; it does not assume that $G-e$ is cubic.
\end{proof}

\begin{corollary}[The remaining girth window for a minimal counterexample]
\label{cor:random-minimal-counterexample-window}
There are absolute constants $K_*,C_3,C_4$ with the following property.
Fix $K\ge K_*$.
Suppose a connected loopless cubic graph $G$ of rank $r$ violates
\[
 c(G)\le2^{r-1}+Kr2^{r/2},
\]
while this bound holds for every connected graph of smaller rank.
If $r$ is sufficiently large, its girth $g$ and minimum cyclic cut order $k$ satisfy
\[
 \frac12\log\log r-\frac{11}{2}\log\log\log r-C_3
 <k=g
 <2\log\log r+2\log\log\log r+C_4.
\]
Moreover, every minimum cyclic cut has a shortest cycle as one shore.
The constants do not depend on $K$.
\end{corollary}

\begin{proof}
Choose $K_*$ at least the absolute constant in the established sharp-order bound for graphs with a cyclic cut of size at most four.
Such cuts are then excluded in $G$.
Corollary~\ref{cor:path-endpoint-loglog-cut} gives the stated lower bound on $k$, since $G$ has cycle count above the half bound.
In particular, $k$ and $g\ge k$ tend uniformly to infinity in the possible counterexamples.

Apply Corollary~\ref{cor:random-improved-girth-window}.
It excludes every girth between
\[
 \left\lceil2\log\log r+2\log\log\log r+C_*\right\rceil
 \quad\text{and}\quad
 \frac12\log r-\frac{11}{2}\log\log r-C_0.
\]
For sufficiently large $r$, the second endpoint exceeds $(\log r)/4$.
Corollary~\ref{cor:random-high-girth-minimal-counterexample}, with $c=1/4$, excludes every girth at least $(\log r)/4$.
The two ranges overlap.
Thus the upper bound on $g$ follows after increasing an absolute constant to account for the ceiling.

Finally, Corollary~\ref{cor:path-high-girth-noncycle-minimum-cut} excludes a minimum cyclic cut with two noncycle shores once $g$ is sufficiently large.
Every minimum cut therefore isolates a cycle.
Its boundary order equals that cycle's length, which is at least $g$.
Since $k\le g$, one has $k=g$, and the isolated cycle is shortest.
\end{proof}
